The MOSI protocol is an extension of the basic MSI cache coherency
protocol. It is a snoop-based protocol. It adds the state Owned~\s{O},
which indicates that the current processor owns this block, and will
service requests from other processors for the block. This also
reduces the amount of write-back data upon cache eviction.

The protocol can be modeled as a parameterized system where the
processes are arranged in a multiset. Each process can communicate
with its neighbors by broadcasting a message on the connecting bus.

\begin{wrapfigure}{r}[2pt]{0.35\linewidth}
  \newcommand{\wrong}{\raisebox{-0.3pt}{\includegraphics[height=1em]{img/skull.png}}}
  \newcommand{\valid}{{\color{green}\bf$\checkmark$}}
  \newcommand{\cache}[1]{\bf\small\emph{#1}}
  \hfill%
  \begin{tabular}{c|cccc}
                & {\cache{M}} & {\cache{O}} & {\cache{S}} & {\cache{I}} \\\hline
    {\cache{M}} & {\wrong}    & {\wrong}    & {\wrong}    & {\valid}       \\
    {\cache{O}} & {\wrong}    & {\wrong}    & {\valid}    & {\valid}       \\
    {\cache{S}} & {\wrong}    & {\valid}    & {\valid}    & {\valid}       \\
    {\cache{I}} & {\valid}    & {\valid}    & {\valid}    & {\valid}       \\
  \end{tabular}
\end{wrapfigure}
%
The state of a cache line can be \emph{Modified}~\s{M},
\emph{Owned}~\s{O}, \emph{Shared}~\s{S}, and \emph{Invalid}~\s{I}.
%
The broadcast communication is depicted using the two following
automata.
%
Initially, all cache lines are in state \emph{Invalid}~\s{I}. The
permitted states of any given pair of cache lines is given in the
table beside.

The first automaton represents the action taken when the given process
issues the broadcast message and/or when it manipulates the cache
line.
%
The cache can be written (\emph{CPUwrite}), caused by a store miss, it
can be read (\emph{CPUread}) caused by a load miss and finally, it can
be replaced (\emph{CPUrepl}).
%
The message sent on the bus vary depending on the state of the current
cache line. The active process can send a \emph{read-to-share}
(\emph{BUSrts}) request, a \emph{read-to-write} (\emph{BUSrtw})
request, a \emph{write-back} (\emph{BUSwb}) request (in case of cache
eviction) or an \emph{invalidate} (\emph{BUSinv}) request. We label
each edge of the automaton with the action
%
\raisebox{-3pt}{\begin{tikzpicture}
  \begin{pgfinterruptboundingbox}
    \path node[message](n){\emph{action}\nodepart{second}\emph{message}} [clip] (n.west) rectangle (n.north east);
  \end{pgfinterruptboundingbox}
  \path[use as bounding box,baseline=(n.one base)] (n.west) rectangle (n.north east);
\end{tikzpicture}}
%
issued by the current process and with the message %
\raisebox{-5pt}{\begin{tikzpicture}
  \begin{pgfinterruptboundingbox}
    \path node[message](n){\emph{action}\nodepart{second}\emph{message}} [clip] (n.south west) rectangle (n.east);
  \end{pgfinterruptboundingbox}
  \path[use as bounding box,baseline=(n.second base)] (n.south west) rectangle (n.east);
\end{tikzpicture}}
%
it sent on the bus. We use
%
\raisebox{-5pt}{\begin{tikzpicture}
  \begin{pgfinterruptboundingbox}
    \path node[message](n){\emph{action}\nodepart{second}$-$} [clip] (n.south west) rectangle (n.east);
  \end{pgfinterruptboundingbox}
  \path[use as bounding box] (n.south west) rectangle (n.east);
\end{tikzpicture}}
%
when no message is sent on the bus.

\smallskip%
{\noindent\centering
  \begin{tikzpicture}[node distance=3cm,
    state/.append style={state-n,inner sep=2pt,scale=1.5},
    message/.append style={scale=0.5}
    ]

    \node[state] (m)	              {M};
    \node[state] (o) [right=4cm of m] {O};
    \node[state] (s) [below=of m]     {S};
    \node[state] (i) [below=of o]     {I};

    \draw [->,myedge] (m) .. controls +(-0.5,1) and +(-1,0.5) ..
    node[message,anchor=north east,pos=0.7]{\emph{CPUread}\nodepart{second}$-$}
    node[message,anchor=south east,pos=0.3]{\emph{CPUwrite}\nodepart{second}$-$} (m);
    \draw [->,myedge] (m) to[out=-30,in=120] node[message,anchor=south,sloped]{\emph{CPUrepl}\nodepart{second}\emph{BUSwb}} (i);
    \draw [->,myedge] (o) -- node[message,above]{\emph{CPUwrite}\nodepart{second}\emph{BUSinv}} (m);
    \draw [->,myedge] (o) -- node[message,right]{\emph{CPUrepl}\nodepart{second}\emph{BUSwb}} (i);
    \draw [->,myedge] (o) .. controls +(0.5,1) and +(1,0.5) .. node[message,anchor=south west]{\emph{CPUread}\nodepart{second}$-$} (o);
    \draw [->,myedge] (s) to[out=-30,in=-150] node[message,anchor=north]{\emph{CPUrepl}\nodepart{second}$-$} (i);
    \draw [->,myedge] (s) .. controls +(-1,-0.5) and +(-0.5,-1) .. node[message,anchor=north east]{\emph{CPUread}\nodepart{second}$-$} (s);
    \draw [->,myedge] (s) -- node[message,left]{\emph{CPUwrite}\nodepart{second}\emph{BUSinv}} (m);
    \draw [->,myedge] (i) -- node[message,anchor=south]{\emph{CPUread}\nodepart{second}\emph{BUSrts}} (s);
    \draw [->,myedge] (i) to[out=150,in=-60] node[message,anchor=north,sloped]{\emph{CPUwrite}\nodepart{second}\emph{BUSrtw}} (m);
    
  \end{tikzpicture}
  \par%
}

\smallskip%
The second automaton represents how other processes react upon
reception of a message %
\raisebox{-5pt}{\begin{tikzpicture}
  \begin{pgfinterruptboundingbox}
    \path node[message](n){\emph{action}\nodepart{second}\emph{message}} [clip] (n.south west) rectangle (n.east);
  \end{pgfinterruptboundingbox}
  \path[use as bounding box,baseline=(n.second base)] (n.south west) rectangle (n.east);
\end{tikzpicture}}.

\smallskip%
{\noindent\centering
  \begin{tikzpicture}[node distance=3cm,
    state/.append style={state-n,inner sep=2pt,scale=1.5},
    message/.append style={scale=0.5,rectangle split parts=1,fill=pink!10!white,rectangle split part fill={yellow!10!white}}
    ]

    \node[state] (m)	              {M};
    \node[state] (o) [right=4cm of m] {O};
    \node[state] (s) [below=of m]     {S};
    \node[state] (i) [below=of o]     {I};
    
    \draw [->,myedge] (m) -- node[message,anchor=south]{\emph{BUSrts} (+Data)} (o);
    \draw [->,myedge] (m) -- node[message,sloped,above]{\emph{BUSrtw} (+Data)} (i);
    \draw [->,myedge] (o) -- node[message,anchor=west] {\ensuremath{\begin{tabular}{l} \emph{BUSrtw} (+Data) \\ \emph{BUSinv}\end{tabular}}} (i);
    \draw [->,myedge] (s) -- node[message,anchor=south east,pos=0.6]{\ensuremath{\begin{tabular}{l}\emph{BUSrtw}\\\emph{BUSinv}\end{tabular}}} (i);
    \draw [->,myedge] (s) .. controls +(-1,-0.5) and +(-0.5,-1) .. node[message,anchor=east]{\ensuremath{\begin{tabular}{l}\emph{BUSrts}\\\emph{BUSwb}\end{tabular}}} (s);
    \draw [->,myedge] (i) .. controls +(1,-0.5) and +(0.5,-1) .. node[message,anchor=south west]{\ensuremath{\begin{tabular}{l}\emph{BUSrts}\\\emph{BUSrtw}\\\emph{BUSinv}\\\emph{BUSwb}\end{tabular}}} (i);
    \draw [->,myedge] (o) .. controls +(0.5,1) and +(1,0.5) .. node[message,anchor=south west]{\emph{BUSrts} (+Data)} (o);
  \end{tikzpicture}
  \par%
}
